\documentclass[10pt,a4paper]{amsart}
\usepackage{amsmath,amssymb,amsthm,graphicx}
\usepackage{verbatim}
\newenvironment{reference}{\comment}{\endcomment}
%\newenvironment{reference}{}{}
\newenvironment{slogan}{\comment}{\endcomment}
\newenvironment{history}{\comment}{\endcomment}

% For commutative diagrams we use Xy-pic
\usepackage[all]{xy}

% We use 2cell for 2-commutative diagrams.
\xyoption{2cell}
% Theorem environments.
%
\theoremstyle{plain}
\newtheorem{theorem}[subsection]{Theorem}
\newtheorem{proposition}[subsection]{Proposition}
\newtheorem{lemma}[subsection]{Lemma}

\theoremstyle{definition}
\newtheorem{definition}[subsection]{Definition}
\newtheorem{example}[subsection]{Example}
\newtheorem{exercise}[subsection]{Exercise}
\newtheorem{situation}[subsection]{Situation}

\theoremstyle{remark}
\newtheorem{remark}[subsection]{Remark}
\newtheorem{remarks}[subsection]{Remarks}

\numberwithin{equation}{subsection}

% Macros
%
\def\lim{\mathop{\mathrm{lim}}\nolimits}
\def\colim{\mathop{\mathrm{colim}}\nolimits}
\def\Spec{\mathop{\mathrm{Spec}}}
\def\Hom{\mathop{\mathrm{Hom}}\nolimits}
\def\Ext{\mathop{\mathrm{Ext}}\nolimits}
\def\SheafHom{\mathop{\mathcal{H}\!\mathit{om}}\nolimits}
\def\SheafExt{\mathop{\mathcal{E}\!\mathit{xt}}\nolimits}
\def\Sch{\mathit{Sch}}
\def\Mor{\mathop{\mathrm{Mor}}\nolimits}
\def\Ob{\mathop{\mathrm{Ob}}\nolimits}
\def\Sh{\mathop{\mathit{Sh}}\nolimits}
\def\NL{\mathop{N\!L}\nolimits}
\def\CH{\mathop{\mathrm{CH}}\nolimits}
\def\proetale{{pro\text{-}\acute{e}tale}}
\def\etale{{\acute{e}tale}}
\def\QCoh{\mathit{QCoh}}
\def\Ker{\mathop{\mathrm{Ker}}}
\def\Im{\mathop{\mathrm{Im}}}
\def\Coker{\mathop{\mathrm{Coker}}}
\def\Coim{\mathop{\mathrm{Coim}}}

% Boxtimes
%
\DeclareMathSymbol{\boxtimes}{\mathbin}{AMSa}{"02}

%
% Macros for moduli stacks/spaces
%
\def\QCohstack{\mathcal{QC}\!\mathit{oh}}
\def\Cohstack{\mathcal{C}\!\mathit{oh}}
\def\Spacesstack{\mathcal{S}\!\mathit{paces}}
\def\Quotfunctor{\mathrm{Quot}}
\def\Hilbfunctor{\mathrm{Hilb}}
\def\Curvesstack{\mathcal{C}\!\mathit{urves}}
\def\Polarizedstack{\mathcal{P}\!\mathit{olarized}}
\def\Complexesstack{\mathcal{C}\!\mathit{omplexes}}
% \Pic is the operator that assigns to X its picard group, usage \Pic(X)
% \Picardstack_{X/B} denotes the Picard stack of X over B
% \Picardfunctor_{X/B} denotes the Picard functor of X over B
\def\Pic{\mathop{\mathrm{Pic}}\nolimits}
\def\Picardstack{\mathcal{P}\!\mathit{ic}}
\def\Picardfunctor{\mathrm{Pic}}
\def\Deformationcategory{\mathcal{D}\!\mathit{ef}}
\title{The sequence of blowups along a nonsingular arc reaches a regular point when one local ring is regular}
\author{The Stacks Project Authors}
\date{stacks-project:resolve.tex:local-snapshot-20261012; GNU FDL 1.2 or later}
\begin{document}
\maketitle
\noindent\textit{Editorial note.} All delivered mathematical text is the author's own, copied from the pinned Stacks Project chapter file resolve.tex; no statement or proof step was rewritten, reordered or removed.\par\medskip
\noindent\textit{Editorial note.} The counted proof is the author's complete argument: the reduction to a minimal generating set of the kernel with the easy cases, the derivation of relations controlled by powers of the uniformizer together with the discarding of generators, and the computation of the transformed relations on the affine blowup algebra showing that the exponents strictly decrease.\par\medskip
\noindent\textit{Editorial note.} The label lemma-sequence-blowups-along-arc-becomes-nonsingular names the counted statement, and the proof environment immediately following it in the source is delivered verbatim as its proof.\par\medskip
\noindent\textit{Editorial note.} The counted statement is carried by the author's own lemma environment in the chapter source.\par\medskip
\noindent\textit{Editorial note.} Pre-open gate: python operations/agent-authoring/tools/precard\_gate.py --source resolve.tex --label lemma-sequence-blowups-along-arc-becomes-nonsingular --start 2896 --end 2954 --local-id 28094 returned RESULT=OK (L28094 free, occupancy FREE, public interval clear). Pre-assembly tag self-check: canonical tag resolved and not present in reference/public-items.jsonl.\par\medskip
\noindent\textit{Editorial note.} The statement and proof are copied from the pinned chapter file; the cited inputs are kept as external pointers with their canonical Stacks tags.\par\medskip
\noindent\textit{Allowed standard prerequisites (editorial).}\begin{itemize}
\item Description of the blowup algebra along the arc. Standard description of the affine blowup algebra of the maximal ideal with respect to an element mapping to a uniformizer, and the extension of the arc to that algebra with kernel the ideal generated by the quotients. Supplies the algebra generated by the quotients of the chosen generators and the new defining equations of the arc.
\item Minimality of the number of generators of the kernel ideal. Standard manipulation of a minimal generating set of an ideal in a local ring, allowing generators to be discarded when a suitable relation becomes trivial. Supplies the induction on the minimal number of generators of the kernel of the surjection onto the discrete valuation ring.
\end{itemize}\medskip
\begin{lemma}
\label{lemma-sequence-blowups-along-arc-becomes-nonsingular}
Let $(A, \mathfrak m, \kappa)$ be a Noetherian local domain of
dimension $2$. Let $A \to R$ be a surjection onto a
complete discrete valuation ring.
This defines a nonsingular arc $a : T = \Spec(R) \to \Spec(A)$. Let
$$
\Spec(A) = X_0 \leftarrow X_1 \leftarrow X_2 \leftarrow X_3 \leftarrow \ldots
$$
be the sequence of blowing ups constructed from $a$.
If $A_\mathfrak p$ is a regular local ring where
$\mathfrak p = \Ker(A \to R)$, then
for some $i$ the scheme $X_i$ is regular at $x_i$.
\end{lemma}

\begin{proof}
Let $x_1 \in \mathfrak m$ map to a uniformizer of $R$.
Observe that $\kappa(\mathfrak p) = K$ is the
fraction field of $R$. Write $\mathfrak p = (x_2, \ldots, x_r)$
with $r$ minimal. If $r = 2$, then $\mathfrak m = (x_1, x_2)$
and $A$ is regular and the lemma is true. Assume $r > 2$.
After renumbering if necessary,
we may assume that $x_2$ maps to a uniformizer of $A_\mathfrak p$.
Then $\mathfrak p/\mathfrak p^2 + (x_2)$ is annihilated by a power
of $x_1$. For $i > 2$ we can find $n_i \geq 0$ and $a_i \in A$
such that
$$
x_1^{n_i} x_i - a_i x_2 = \sum\nolimits_{2 \leq j \leq k} a_{jk} x_jx_k
$$
for some $a_{jk} \in A$. If $n_i = 0$ for some $i$, then we can remove
$x_i$ from the list of generators of $\mathfrak p$ and we win by
induction on $r$. If for some $i$ the element $a_i$ is a unit, then
we can remove $x_2$ from the list of generators of $\mathfrak p$
and we win in the same manner. Thus either
$a_i \in \mathfrak p$ or $a_i = u_i x_1^{m_1} \bmod \mathfrak p$
for some $m_1 > 0$ and unit $u_i \in A$. Thus we have either
$$
x_1^{n_i} x_i = \sum\nolimits_{2 \leq j \leq k} a_{jk} x_jx_k
\quad\text{or}\quad
x_1^{n_i} x_i - u_i x_1^{m_i} x_2 =
\sum\nolimits_{2 \leq j \leq k} a_{jk} x_jx_k
$$
We will prove that after blowing up the integers $n_i$, $m_i$
decrease which will finish the proof.

\medskip\noindent
Let us see what happens with these equations on the affine blowup
algebra $A' = A[\mathfrak m/x_1]$. As $\mathfrak m = (x_1, \ldots, x_r)$
we see that $A'$ is generated over $R$ by $y_i = x_i/x_1$ for $i \geq 2$.
Clearly $A \to R$ extends to $A' \to R$ with kernel
$(y_2, \ldots, y_r)$. Then we see that either
$$
x_1^{n_i - 1} y_i = \sum\nolimits_{2 \leq j \leq k} a_{jk} y_jy_k
\quad\text{or}\quad
x_1^{n_i - 1} y_i - u_i x_1^{m_1 - 1} y_2 =
\sum\nolimits_{2 \leq j \leq k} a_{jk} y_jy_k
$$
and the proof is complete.
\end{proof}

\end{document}
